\documentclass[10pt]{article}
\usepackage[margin=0.72in]{geometry}
\usepackage{booktabs}
\usepackage{microtype}
\usepackage[hidelinks]{hyperref}
\usepackage{titlesec}
\usepackage{parskip}

\titlespacing*{\section}{0pt}{6pt}{3pt}
\titlespacing*{\subsection}{0pt}{5pt}{2pt}
\titleformat{\section}{\large\bfseries}{\thesection.}{0.5em}{}
\titleformat{\subsection}{\normalsize\bfseries}{}{0em}{}
\setlength{\parskip}{3pt}
\pagestyle{empty}
\newcommand{\code}[1]{\texttt{\small #1}}

\begin{document}

\begin{center}
{\Large\bfseries The Assert Line Is Doing the Work:\\[1pt]
An Ablation of EvaLoop's Carried Example}\\[3pt]
{\small CSC 591/791 --- Trustworthy AI for Software Engineering, Assignment 4}
\end{center}

\vspace{-3pt}

\section{Extension and Research Question}

EvaLoop measures robustness by feeding a model its own output: generate code
from a description, test it, summarise the code back to a description, repeat.
One detail keeps the loop runnable. Before a summary becomes the next cycle's
prompt it is wrapped as a docstring \emph{with the final assert line of the
original prompt appended}, so the function name and signature survive. The
paper presents this as plumbing. But that line is not only a signature --- it is
a \textbf{worked example}: one concrete input with its expected output,
re-injected from outside the loop into every single cycle. The model never has
to recover the task's interface from its own summary, and it always sees one
ground-truth test case.

\textbf{Research question: how much of EvaLoop's measured robustness rests on
that leaked example?} I added a new loop type that keeps the standard loop
intact and weakens only what the summary carries forward, in four arms that
remove information monotonically:

\begin{center}\small
\begin{tabular}{llc}
\toprule
arm & carrier appended to the summary & information retained \\
\midrule
\code{assert} & \code{assert f(3, 4) == 7} & name, arity, example \\
\code{signature} & \code{def f(arg1, arg2):} & name, arity \\
\code{name} & \code{The function must be named `f`.} & name \\
\code{bare} & \emph{(nothing)} & nothing \\
\bottomrule
\end{tabular}
\end{center}

The \code{assert} arm is a re-implementation of the standard loop and serves as
my own baseline; differences between adjacent arms attribute the standard
loop's score to the example, the arity and the name respectively. Signature
stability is solved rather than dodged: the \code{signature} and \code{name}
arms still pin the function name, so a failure means the model could not
re-derive the task from its own description, not that it renamed the function.
The \code{bare} arm gives that up deliberately and is reported as the floor.

\section{The Framework and My Changes}

A cycle is a list of tasks chained by \code{BatchEvaluationCycle}: the output of
task $n$ is the input of task $n{+}1$, and the output of the task flagged
\code{produces\_testable\_code} is run against the EvalPlus MBPP Plus suites
(\emph{base} = original tests, \emph{plus} = extended; a cycle passes only if
both do). A task implements three hooks --- \code{create\_prompt},
\code{extract\_output}, \code{postprocess} --- and registers itself with
\code{@register\_task}. The assert re-attachment lives in
\code{CodeSummarizationTask.postprocess}, which recovers the line as
\code{context["cycle\_input"].splitlines()[-2]}. ASL is
$\sum_i n_i i^2 s_i / (TM)$ in \code{analysis/metrics.py}.

So a new loop type is almost entirely new code. I added two files:
\code{tasks/assert\_ablation.py} (the four arms plus the signature recovery) and
\code{models/gemini\_vertex.py} (a Gemini/Vertex backend --- the framework had
none, and Gemini ids fell through to the HuggingFace branch). Six existing
files changed: \code{tasks/\_\_init\_\_.py} (register), \code{core/config.py}
(a \code{loop\_type} field, the arms in \code{get\_task\_configs}, Gemini
routing, \code{max\_tasks}), \code{core/evaluator.py} (\code{loop\_type} was
hard-coded; result files now carry the arm and a \code{run\_metadata} block
with measured token usage), \code{models/factory.py} (register Gemini; make the
torch/vLLM imports lazy so an API-only run does not need a 2.5\,GB install),
and \code{cli.py} (the new flags).

\textbf{The carrier is derived once, from the original prompt.}
\code{signature\_from\_assert} parses the cycle-1 assert line with
\code{ast} --- picking the leftmost non-builtin call, since the task function is
often wrapped (\code{assert set(similar\_elements(...)) == ...}) and a few tasks
are named after a builtin --- and the result is memoised per \code{task\_id}.
Later cycles reuse the memo instead of re-reading their own docstring, so every
arm's carrier traces to the same source as the baseline's and cannot itself
drift. Validated against all 378 MBPP Plus assert lines: it recovers the entry
point for 378/378.

\section{Experimental Setup}

Two models reachable through the course's Vertex credits,
\code{gemini-2.5-flash} and \code{gemini-3.5-flash-lite}; greedy decoding
($T{=}0$, top-$p$ 1.0), 1024 new tokens, max 10 cycles, one run per
configuration; 8 runs in all. Tasks are the first 120 of EvalPlus's MBPP Plus
ordering (the full 378 was beyond the project's request quota in the available
time; the subset is a deterministic prefix and the \code{--max\_tasks} flag
reproduces it). Measured cost: \textbf{\$3.08} over 14{,}342 calls (1.54M input
/ 1.05M output tokens), \$0.0032 per task-arm.

Two setup findings are worth recording because both would have silently
corrupted the results.

\textbf{Reasoning tokens starve the answer.} Gemini 2.5+ count reasoning
against \code{max\_output\_tokens}. At EvaLoop's 1024-token budget, 35 of 120
cycle-1 generations were cut off mid-expression (\code{for i in range(1, n}),
pushing pass@1 to 62.5\% against a published 80.0\% on the same 120 tasks. One
call consumed 1020 output tokens to emit 161 characters. Since the protocol
specifies 1024 tokens \emph{for the answer}, I disable thinking; the same call
then uses 65 output tokens and completes.

\textbf{Abandoned API calls look like model failures.} An exhausted retry
returns an empty generation, which the loop records as a broken cycle. Running
four arms concurrently exhausted retries on ${\sim}1\%$ of calls --- unevenly
across arms, i.e.\ exactly on the quantity being measured. The backend therefore
paces request starts with an AIMD limiter and counts abandoned calls; a run
with a non-zero count is marked untrustworthy. All 8 reported runs have
\textbf{\code{hard\_failures}~=~0}.

\section{Results and Failure Analysis}

\begin{center}\small
\begin{tabular}{llcccc}
\toprule
model & arm & pass@1 & ASL\_pow & vs \code{assert} & mid-loop failures \\
\midrule
\code{gemini-2.5-flash} & \code{assert} & 85.0\% & 8.115 & --- & 5 \\
 & \code{signature} & 83.3\% & 7.974 & $-1.7\%$ & 6 \\
 & \code{name} & 82.5\% & 7.809 & $-3.8\%$ & 6 \\
 & \code{bare} & 81.7\% & 6.593 & $-18.8\%$ & 21 \\
\midrule
\code{gemini-3.5-flash-lite} & \code{assert} & 85.0\% & 7.982 & --- & 7 \\
 & \code{signature} & 83.3\% & 7.324 & $-8.2\%$ & 15 \\
 & \code{name} & 84.2\% & 7.705 & $-3.5\%$ & 10 \\
 & \code{bare} & 83.3\% & \textbf{1.794} & $\mathbf{-77.5\%}$ & 85 \\
\bottomrule
\end{tabular}
\end{center}

\textbf{Internal validity.} pass@1 is flat across arms (81.7--85.0\%) while
mid-loop failures move by more than an order of magnitude. That is what the
design predicts: cycle 1's input is the untouched prompt in every arm, so the
arms must agree there and may only diverge from cycle 2, where the carrier
enters. The protocol check puts my \code{assert} arm at ASL\_pow 8.115 against
the published 6.923 (all 378 tasks) and 7.167 (the same 120). Being modestly
above is consistent with disabling thinking and with model drift since the data
was collected; the absolute values are therefore not interchangeable with the
leaderboard's, and I rely only on within-experiment comparisons.

\textbf{The assert line compresses the ranking.} Under the standard loop the
two models look equally robust: 8.115 versus 7.982, 1.7\% apart. Remove the
carrier and they are 6.593 versus 1.794 --- a factor of 3.7. The leaked example
is not a small convenience; it \emph{hides} a large difference in how well each
model can reconstruct a task from its own prose. \code{gemini-2.5-flash}
degrades gracefully (77 of 120 tasks still survive 10 cycles bare);
\code{gemini-3.5-flash-lite} collapses (15 survive, 49 break at cycle 2).

\textbf{Failure modes.} I adjudicated 35 failed loops --- 4 per arm plus 3 from
the provided baseline, stratified by arm, excluding loops that broke at cycle 1
because the arms are identical there by construction. The causes partition
almost exactly along what each arm removed:

\begin{center}\small
\begin{tabular}{lcccccc}
\toprule
arm & signature-drift & lost-example & dropped-req. & cannot-impl. & plus-only & other \\
\midrule
\code{assert} & 0 & 0 & 4 & 0 & 4 & 0 \\
\code{signature} & 0 & 4 & 0 & 0 & 4 & 0 \\
\code{name} & 0 & 2 & 2 & 1 & 3 & 0 \\
\code{bare} & \textbf{6} & 2 & 0 & 0 & 0 & 0 \\
provided baseline & 0 & 0 & 0 & 0 & 0 & \textbf{3} \\
\bottomrule
\end{tabular}
\end{center}

\emph{A loop the extension broke.} \code{Mbpp/69} (\code{is\_sublist}) fails in
three different arms for the same reason and never in \code{assert}. The
summary is faithful --- ``determines if \code{sub\_list} is a contiguous
sequence within \code{main\_list}'' --- but the test calls
\code{is\_sublist(main, sub)}. The original assert line shows that order; a
prose summary does not. Without it the model writes
\code{def is\_sublist(sub\_list, main\_list)} and fails on argument order, not
on semantics. This is the mechanism the ablation was built to expose, and it is
invisible to the standard loop.

\emph{A loop it did not break.} \code{Mbpp/244} (\code{next\_Perfect\_Square})
fails identically in \code{assert} and \code{signature} at cycle 1 with
\code{base} passing and only \code{plus} failing: an EvalPlus edge case, not
drift. Such \code{plus-only} failures are uniform background noise across arms
(4, 4, 3, 0) and the carrier has no bearing on them.

\emph{The provided baseline contains infrastructure failures.} All three
sampled baseline failures have \code{Error generating response: 'NoneType'
object is not subscriptable} stored where the generated code belongs --- API
errors recorded as loop failures --- and the LLM judge scored each
\code{boundary\_similarity}~=~1.0, so the published pipeline's own drift check
did not flag them. This is the contamination my \code{hard\_failures} counter
exists to prevent, and it means the published ASL values are, to some unknown
extent, measuring API reliability.

\section{Limitations and Implications}

\textbf{Two models is a pairwise comparison, not a ranking.} I can show how
much each model depends on the carrier; I cannot claim the ablation reorders a
leaderboard. The ordering in fact does not flip --- what changes is the
magnitude, by $11\times$.

\textbf{The gradient is not monotone, and that is informative.} For
\code{gemini-3.5-flash-lite}, \code{signature} ($-8.2\%$) hurts \emph{more}
than \code{name} ($-3.5\%$), although it carries strictly more information.
Inspection suggests \code{def f(arg1, arg2):} actively misleads: the model
adopts the meaningless parameter names and loses the semantic cues it would
otherwise have taken from the description. A carrier can be worse than a weaker
carrier if it displaces better information --- so ``more context'' is not
monotonically safer.

\textbf{What the numbers do not establish.} 120 tasks leaves roughly 7 mid-loop
failures per well-behaved arm, so the \code{assert}/\code{signature}/\code{name}
differences (1.7--8.2\%) rest on single-digit counts and I would not defend
their ordering; only the \code{bare} effect ($-18.8\%$, $-77.5\%$) is large
relative to that noise. One run per configuration means no variance estimate.
Thinking is disabled, which is a deviation from whatever configuration produced
the published data. Finally, the 35 failure labels were produced by an LLM
reading each loop's full source context and then reviewed and accepted by the
author, not assigned by independent annotators; the worksheet carrying the
evidence for every label ships with the submission so they can be rechecked.

\textbf{Implication.} EvaLoop's headline claim is that self-consistency
measures robustness without external perturbations. The ablation shows that the
standard loop also quietly supplies the model with a ground-truth test case
every cycle, and that this single leak accounts for most of the measured
robustness of a weaker model --- 77.5\% of \code{gemini-3.5-flash-lite}'s ASL.
A benchmark whose scores are produced partly by information it injects itself
will compress the differences it exists to reveal. If the goal is to measure
whether a model can sustain its own description of a task, the interface should
be carried by something that cannot double as an answer --- a signature derived
from the model's own code, for instance --- and the carrier should be reported
as a parameter of the benchmark, not left as plumbing.

\end{document}
